\chapter{Quick Start}
\label{ch:geolog-quickstart}

In this chapter you will install geolog, write a theory and a handful
of facts in one file, and run the chase: the engine will derive the
tables your axioms entail, invent nothing you did not license, and
report --- rather than repair or reject --- the one state your theory
forbids. At the end you will look at the theory you wrote and find it
is itself rows. The whole path fits in one sitting; every command is
copy-pasteable, and every output shown is the verbatim transcript of
a real run.\footnote{The book's build regenerates these transcripts
by executing the printed commands against the committed examples,
and each fragment's header records the exact tool revision it ran
against (all runs on one commodity arm64 box).}

\section{Install}
\label{sec:gqs-install}

A Rust toolchain is the only prerequisite (stable channel;
\path{rust-toolchain.toml} pins the components).

\begin{Verbatim}[fontsize=\small,frame=leftline,rulecolor=\color{accentmid},framesep=3mm,xleftmargin=3mm]
git clone https://github.com/maitria-coalition/geolog-alpha
cd geolog-alpha
cargo install --path crates/geolog
\end{Verbatim}

\noindent A fresh clone compiled in about a minute on the commodity
arm64 box that regenerates this book's transcripts.
\texttt{geolog --help} lists ten verbs; the three
this chapter uses are \texttt{check}, \texttt{chase}, and
\texttt{rows}.

\section{A theory}
\label{sec:gqs-theory}

The example is a freight network, small enough to hold in your head:
places, direct hops, and reachability. The file
\path{examples/routes.geolog} declares one sort, two relations, and
three axioms:

\begin{Verbatim}[fontsize=\small,frame=leftline,rulecolor=\color{accentmid},framesep=3mm,xleftmargin=3mm]
theory Routes {

  Node : Set;

  // Asserted: a direct hop exists.
  Edge : [src: Node, tgt: Node] -> Prop;

  // Derived: some chain of hops exists. Rows appear exactly when the
  // axioms force them — never asserted directly.
  Path : [src: Node, tgt: Node] -> Prop;

  // `chase` marks a derivation rule: the axiom FIRES, creating rows.
  // (`report` marks a constraint: violations are judged, nothing is
  // created; an unmarked axiom defaults to report — the least-
  // dangerous reading, so derivation is always an explicit choice.)
  chase ax/path/edge : forall a : Node, b : Node.
    [src: a, tgt: b] Edge |- [src: a, tgt: b] Path;

  chase ax/path/trans : forall a : Node, b : Node, c : Node.
    [src: a, tgt: b] Path, [src: b, tgt: c] Edge
    |- [src: a, tgt: c] Path;

  // A denial: this network is required loop-free. A commit that lands
  // a cycle is not rejected — the violated denial is REPORTED as a
  // judgement, and repair is a further ordinary commit.
  report ax/no/loop : forall a : Node.
    [src: a, tgt: a] Path |- false;
}
\end{Verbatim}


Four things to read off the shape.
\begin{enumerate}
\item Relations are declared over named record fields, and axioms
are \emph{sequents}: universally quantified entailments, each
premise a conjunction of atoms, each conclusion here a single atom.
The odd one out is the third axiom, whose conclusion \texttt{false}
makes it a \emph{denial}: a state the theory forbids.
\item The word in front of each axiom's name is its \emph{mode
mark}, and the file's own comments state the semantics:
\texttt{chase} fires, \texttt{report} judges, and an unmarked axiom
means \texttt{report}. An axiom creates rows only where its author
said so in as many words.
\item Note what the mark does \emph{not} do: nothing in the
\emph{declarations} marks \texttt{Path} as derived. The mode belongs
to the axiom, and \texttt{Path} being a table nobody writes remains
an emergent fact --- about which relations your \texttt{chase}-marked
axioms conclude into and which your commits assert into --- not a
schema annotation.
\item The slash in \texttt{ax/path/edge} is nothing special:
slash-names are ordinary identifiers, used by convention for
namespacing. The \texttt{ax/} prefix, however, does real grammatical
work: it makes a statement an axiom rather than a declaration, so a
sequent written under a name like \texttt{my/rule} is a syntax
error. The Reference Manual (chapter~\ref{ch:geolog-reference})
spells out these name rules.
\end{enumerate}

Nothing travels beside the theory: a \path{.geolog} file is
self-contained --- everything it draws on is either declared in it
or imported at the top. This is the \emph{roster-totality}
discipline: every declared relation has a decided evaluator.
\texttt{ordinary} (asserted or chase-derived) is the default,
declared by writing nothing. When a relation is instead decided by a
bounded external procedure, an import line says so ---
\texttt{import <provider>.\{rel\};} --- classifying it as a
\emph{side-condition} awaiting that provider. The same import syntax
authorizes builtin ground vocabulary by name
(\texttt{import std.\{Q, q\_lt\};}). Nothing guesses on
your behalf: an \texttt{import std} line naming something outside
the builtin roster refuses at once, listing the importable names.
This chapter's example needs no imports; every relation in it is
ordinary, and the file's opening comment says so.

\section{An instance}
\label{sec:gqs-instance}

The same file carries the facts, in an \texttt{instance} block ---
element declarations first, then facts in the same record idiom the
axioms use:

\begin{Verbatim}[fontsize=\small,frame=leftline,rulecolor=\color{accentmid},framesep=3mm,xleftmargin=3mm]
instance Demo : Routes {
  hub, east, west, port : Node;

  [src: hub,  tgt: east] Edge;
  [src: hub,  tgt: west] Edge;
  [src: east, tgt: port] Edge;
}
\end{Verbatim}


\noindent Elements are declared before use; a fact naming an
undeclared element is an error, never a silently minted row.
(\texttt{geolog check} validates all of this --- declarations, field
sets, sorts --- without running anything.)

The map is worth ten seconds: \texttt{hub} ships to \texttt{east}
and to \texttt{west}, and \texttt{east} ships on to \texttt{port}.
Before running anything, guess which pairs the axioms will put in
\texttt{Path}.

\section{The chase}
\label{sec:gqs-chase}

\begin{Verbatim}[fontsize=\small,frame=leftline,rulecolor=\color{accentmid},framesep=3mm,xleftmargin=3mm]
$ geolog chase examples/routes.geolog
sorts:
  Node — 4 elements (4 seeded)
relations:
  Edge — 3 rows (3 seeded, 0 derived)
    [src: east, tgt: port] Edge
    [src: hub, tgt: east] Edge
    [src: hub, tgt: west] Edge
  Path — 4 rows (0 seeded, 4 derived)
  + [src: east, tgt: port] Path
  + [src: hub, tgt: east] Path
  + [src: hub, tgt: port] Path
  + [src: hub, tgt: west] Path
fixpoint: reached, no judgements — the model is saturated.
$ echo $?
0
\end{Verbatim}


Look at the \texttt{+} rows: tables nobody wrote.
\texttt{[src: hub, tgt: port]} is in \texttt{Path} because
\texttt{ax/path/trans} forces it from two rows you did write, and the
chase ran that forcing to its \emph{least fixpoint} --- everything the
axioms entail, nothing they do not. (The other three are your own
hops, lifted into \texttt{Path} by \texttt{ax/path/edge}: one new
reachability fact, four new rows.) ``Derived'' names provenance
rather than any lesser grade of truth: these rows are as queryable
as the seeded ones, and rerunning from the same facts re-derives
exactly the same table.

\section{Break it}
\label{sec:gqs-break}

Add one hop to the instance block --- the return leg that closes a
loop --- and save the edited file as \path{routes-cycle.geolog}:

\begin{Verbatim}[fontsize=\small,frame=leftline,rulecolor=\color{accentmid},framesep=3mm,xleftmargin=3mm]
  [src: port, tgt: hub]  Edge;
\end{Verbatim}

\begin{Verbatim}[fontsize=\small,frame=leftline,rulecolor=\color{accentmid},framesep=3mm,xleftmargin=3mm]
$ geolog chase routes-cycle.geolog
sorts:
  Node — 4 elements (4 seeded)
relations:
  Edge — 4 rows (4 seeded, 0 derived)
    ⋮
  Path — 12 rows (0 seeded, 12 derived)
    ⋮
judgements (3):
  denial violated: ax/no/loop at [a: hub]
  denial violated: ax/no/loop at [a: east]
  denial violated: ax/no/loop at [a: port]
$ echo $?
1
\end{Verbatim}


Read what did \emph{not} happen. The write was not rejected: the
facts landed, and the chase saturated over them (twelve \texttt{Path}
rows now: each of the three places on the loop reaches all four,
itself included). The
violation was not repaired: no engine deleted your hop or picked a
row to sacrifice. What happened is a \emph{judgment} --- the very
thing \texttt{ax/no/loop}'s \texttt{report} mark asks for: the
violated denial is reported, with the axiom's own name and the
binding that witnesses it, and the exit code says so. Three
judgments here, not one: the same loop, witnessed once from each
place standing on it. If you have
operated an event-sourced system, this is a familiar posture done
formally: an
append-only log can land in a state some invariant marks invalid, and
the repair is a further ordinary commit --- made by whoever owns the
workflow, not by the database guessing. The Tutorial
(chapter~\ref{ch:geolog-tutorial}) makes this the centre of the
story.

\section{The theory is rows}
\label{sec:gqs-rows}

One last look before you go:

\begin{Verbatim}[fontsize=\small,frame=leftline,rulecolor=\color{accentmid},framesep=3mm,xleftmargin=3mm]
$ geolog rows examples/routes.geolog
// Declaration rows in the GeologMeta idiom (the schema you wrote is
// itself data). refs: L<n> = local element minted at op[n]; G<name> =
// genesis identity. This is the declaration slice; axioms are rows
// too, in the same substrate — the committed listings under
// theories/generated/ carry complete worked examples.
op[   0] elta sort=G<Theory>                    // Theory "Routes"
op[   1] elta sort=G<Sort>                      // Sort "Node"
op[   2] funa fun=G<Sort/name> args=[L1] value=str:"Node"
op[   3] funa fun=G<Sort/theory> args=[L1] value=L0
op[   4] elta sort=G<Fun>                       // Fun "Edge"
op[   5] funa fun=G<Fun/name> args=[L4] value=str:"Edge"
    ⋮
op[   8] elta sort=G<Field>                     // Field "src" of Edge
op[   9] funa fun=G<Field/fun> args=[L8] value=L4
    ⋮
\end{Verbatim}


The theory you wrote in \S\ref{sec:gqs-theory} is itself data: a
\texttt{Theory} element, a \texttt{Sort} element with a name row and
a home row, \texttt{Fun} and \texttt{Field} elements for each
relation and its columns --- rows in the \emph{meta-schema}, the
schema language that is itself a geolog theory. As the listing's own
header says, this is the declaration slice; axioms continue as rows
in the same idiom, and the repository's \path{theories/generated/}
directory carries complete committed listings --- the full Petri-net
theory runs to 1{,}239 ops.
Schemas, extensions, and migrations live in the same substrate they
govern, so you can version them, diff them, and derive facts about
them just as you would any other rows.

\section{Where to next}
\label{sec:gqs-next}

The Tutorial (chapter~\ref{ch:geolog-tutorial}) works an event-sourced
order flow with two writers: views only the merged logs can derive,
two replicas' records identified by congruence rather than a conflict
resolver, a witness the chase invents under an existential axiom, a
violation that only the merge surfaces, and the signed commit log
underneath, which can say things a broker cannot. The Reference Manual
(chapter~\ref{ch:geolog-reference}) gives the systematic account. The
part closes with diagrams as data (chapter~\ref{ch:qtsl-scomplexes}):
one incidence genus, stored and checked as rows, for every diagram
the wider toolkit traffics in. And
for the road from a model to a checked temporal claim about it ---
the reason this database exists in this book --- start the QTSL part
at chapter~\ref{ch:qtsl-quickstart}.

